% Seção: Conclusão  (orçamento: 0,5 página)
%
% Responde às três perguntas de pesquisa da Introdução e aponta direções futuras ligadas às limitações da
% Discussão. Sem frase de fechamento.

\section{Conclusão}\label{sec:conclusão}

Este trabalho comparou \nArms{} abordagens de aprendizado em contexto para responder a perguntas sobre os eventos
da final da Copa do Mundo de 2022, em \nModels{} modelos. Nenhuma abordagem acertou mais em todos os modelos, mas a
abordagem \armGraphTools{} foi a única que manteve o acerto alto em todos, de \accRangeGraphTools{}. A \armTextToCypher{} e o
\armEventsInPrompt{} variaram entre os modelos e ficaram abaixo dela nos menores. Nos maiores, com raciocínio
estendido, a \armTextToCypher{} chegou ao mesmo nível, e o \armEventsInPrompt{} se aproximou, este ao custo de mais
tokens de saída e mais tempo por pergunta. A \armVector{} e as \armStatsInPrompt{} acertaram pouco em todos os modelos.

Os resultados vêm de uma partida e de uma repetição por pergunta. Repetir as execuções permitiria medir a variação
entre rodadas e testar as diferenças entre abordagens. Outros caminhos são estender as perguntas a várias partidas;
avaliar modelos maiores sem raciocínio estendido, para separar o efeito do tamanho do efeito do raciocínio; aplicar
as mesmas ferramentas a um banco relacional (SQL), sem o grafo, para separar o efeito das ferramentas do efeito do
grafo; e testar uma busca vetorial em que cada trecho traga também os eventos vizinhos, como em Tzikas et
al.~\cite{tzikas2025dynamic}. Por fim, as mesmas abordagens podem ser testadas em outros registros de eventos, como
os de processos industriais, com a preparação dos dados, o grafo e as ferramentas feitos para as perguntas desse
domínio.